\section{Controller Design}
\label{appendix:implementation}

All controllers inherit from a common \texttt{BaseController} interface,
accept state $s = [x,\,\dot{x},\,\theta,\,\dot{\theta}]^\top$, and
return a cart-velocity set-point $u$ in m/s. The set-point is
saturated at the axis limit ($\pm 0.9$\,m/s) by the shared command
interface: the simulator clips it before the motor model, and the
hardware bridge applies the identical clip before converting to mm/s.
Saturation is therefore a property of the plant interface rather than
of any individual controller, and all controller families command
against the same effective actuation range in simulation and on
hardware.
This command is passed directly to \texttt{MC\_MoveVelocity}, the
PLCopen velocity-command function block accepted by the Beckhoff servo
drive (it starts an endless travel at a commanded velocity/direction,
not a jogging primitive). Internally, the
servo executes a 7-phase optimized S-curve velocity profile; for control
design we approximate this as a first-order lag ($\tau_v = 0.15$\,s).
Learning-based controllers (RL, MPPI, IQL) are implemented in
PyTorch~\cite{paszke2019pytorch}.
The justification is pragmatic, not algorithmic: (1)~the first-order
lag was identified empirically from open-loop trajectory matching and
captures the dominant outer-loop velocity response that every controller
family actually commands against; and (2)~none of our controllers
(feedback laws, planners, or RL policies) issue jerk/acceleration
set-points, so the S-curve's internal jerk-limited transitions are not a
design degree of freedom for them. Where the higher-order motor
response matters (predictive controllers trying to brake into a wall),
the residual mismatch is characterized as a sim-to-real gap in
Appendix~\ref{appendix:simulation_validation} rather than modeled in
the loop.
This section describes the key design choices for each controller;
selected deployed parameters are listed in Table~\ref{tab:controller_params}.

\subsection{Tuning Methodology}
\label{sec:tuning}

We tune the classical controllers via CMA-ES~\cite{hansen2016cma} (pycma~\cite{hansen2019pycma})
(Covariance Matrix Adaptation Evolution Strategy: a gradient-free
optimizer that adapts a multivariate Gaussian search distribution,
well-suited for low-dimensional, non-convex parameter optimization)
in a first-order motor simulation (RK4, 20~Hz) using an ISE-based objective:
\begin{equation}
  J = 500 \cdot \mathrm{ISE}_\theta + 20 \cdot \mathrm{ISE}_{\dot\theta}
      + 0.5 \cdot \mathrm{ISE}_u + 0.1 \cdot N_\text{boundary}
\end{equation}
where $\mathrm{ISE}_x = \tfrac{1}{T}\int_0^T x(t)^2\,dt$ is the
time-averaged squared error of the ball angle $\theta$ and the
command $u$; the velocity term additionally weights $\dot\theta^2$ by a
Gaussian gate $\exp(-\theta^2/2\sigma_g^2)$ with $\sigma_g \approx 0.0085$
(the value \texttt{0.020/2.355} used in code, applied to the ball
angle $\theta$), so angular velocity is penalized mainly near
equilibrium; and $N_\text{boundary}$ counts
timesteps at which the ball angle reaches its physical limit
($|\theta| \geq \theta_\text{max}$). $J$ is averaged over 7
representative initial conditions spanning all strata.

\subsection{PD Controller}

We use a PD controller whose derivative term acts on the measured angular
velocity:
\begin{equation}
  u_t = K_p\,\theta_t + K_d\,\dot{\theta}_t.
  \label{eq:pd_law}
\end{equation}
No integral term is used ($K_i = 0$, confirmed by CMA-ES convergence):
adding integral action causes windup during wall-override intervals
(when the nominal command is replaced) and was consistently rejected by
the optimizer. The arc's concave dip is passively restoring about the
equilibrium, so no integral term is needed to remove steady-state
angle error in the linear regime.

\paragraph{Wall-proximity override.}
Let $u_\text{nom}$ be the controller's nominal command. When the cart is
within $d_w$ of a rail limit and $u_\text{nom}$ pushes toward the wall, the
output is overridden with a fixed escape velocity $u_\text{ov}$ (m/s):
\begin{equation}
  u_t = \begin{cases}
    -u_\text{ov} & x > x_\text{lim} - d_w \text{ and } u_\text{nom} > 0,\\
    +u_\text{ov} & x < -(x_\text{lim} - d_w) \text{ and } u_\text{nom} < 0,\\
    \ \, u_\text{nom} & \text{otherwise.}
  \end{cases}
  \label{eq:override}
\end{equation}
With the override and calibrated sensors, PD+WO reaches 96\% hardware SR.
Two ablations isolate the contributions: removing the override drops it to
18\% (vanilla PD), and keeping the override but disabling sensor-offset
correction (Appendix~\ref{appendix:sensor_offset}) drops it to 88\%
(Appendix~\ref{appendix:results}). The override prevents irreversible wall
collisions. The same override is applied to LQR.

\subsection{Linear Quadratic Regulator (LQR)}

LQR linearizes the first-order velocity-input dynamics
about $\theta = 0$; the discrete-time algebraic
Riccati equation (DARE)~\cite{lewis2012optimal} yields a state-feedback
gain $K = [K_x,\, K_{\dot x},\, K_\theta,\, K_{\dot\theta}]$. The deployed
wall-override configuration does not use this DARE solution: on hardware the
gain that transfers is the PD-equivalent one (cart terms zero,
Table~\ref{tab:controller_params}), as the DARE and CMA-ES solutions perform
far worse on the real system (see the LQR discussion in the classical-tuning
appendix).
Three variants are evaluated:
\begin{itemize}
  \item \textbf{Basic}: No constraint handling; same PD-equivalent
    deployed gain with the override disabled (20\% HW SR).
  \item \textbf{Q-center}: High $Q_x$ penalizes cart displacement
        (28\% HW SR, oscillatory).
  \item \textbf{Wall override}: Eq.~\eqref{eq:override} with
        $d_w{=}0.12$~m, $u_\text{ov}{=}0.5$ (\textbf{92\%} HW SR, PD-equivalent gain).
\end{itemize}

\subsection{Sliding Mode Controller (SMC)}
\label{sec:smc_impl}

Sliding mode control has been applied to the ball-and-arc system in prior
work~\cite{ho2010smc,kao2018secondorder}. Our implementation defines a
sliding surface $\sigma = \dot\theta + \lambda\theta$
and computes the equivalent control via model inversion~\cite{slotine1991applied}:
\begin{equation}
  v_\text{eq} = -\frac{F}{G}, \quad
  F = f_\text{vel}(\theta,\dot\theta,\dot{x}) + \lambda\dot\theta,\quad
  G = g_\text{vel}(\theta)
  \label{eq:smc}
\end{equation}
with a smooth reaching law $v_\text{reach} = -(k/G)\,\tanh(\sigma/\phi)$
and an explicit cart-velocity damping term:
\begin{equation}
  v_\text{cmd} = -\frac{F + k\,\tanh(\sigma/\phi)}{G} - k_d\,\dot{x}.
  \label{eq:smc_full}
\end{equation}

\paragraph{Cart-velocity damping.}
The term $-k_d\dot{x}$ is a separate cart-velocity damper appended to
the SMC command; it is not part of the sliding surface (which depends
only on $\dot\theta$ and $\theta$). The ball bearing's rolling friction
(Appendix~\ref{appendix:dynamics}) enters $f_\text{vel}$ and is included in the
nominal equivalent-control compensation; it is not what $k_d$ addresses. Instead, $k_d$
closes an open-loop cart-drift mode: when $\sigma\approx 0$ and
$\theta\approx 0$ the model-inversion terms in $v_\text{eq}$ collapse to
small residuals (model error, discretization, sensor offset) that act as
a constant velocity command; without damping these integrate through
the first-order cart lag into unbounded cart motion after the ball is
balanced. Adding $-k_d\dot{x}$ directly opposes any residual cart
velocity, independent of ball angle.  We find $k_d{=}0.5$ sufficient
on hardware.

\paragraph{Boundary layer.}
The boundary layer thickness $\phi$ controls the trade-off between
chattering and steady-state accuracy.  Small $\phi$ (${<}0.15$)
causes chattering ($>$500 sign reversals per episode at 20~Hz);
large $\phi$ (${>}0.8$) reduces $\tanh$ to a near-linear function,
effectively converting the controller to a continuous nonlinear
approximation rather than a true variable-structure controller.
CMA-ES selects $\phi{=}0.806$ for the deployed SMC+WO variant,
in the large-$\phi$ regime. We acknowledge this trade-off:
the sliding-mode discontinuity is sacrificed for
stability on hardware at 20~Hz with noisy sensor feedback;
smaller $\phi$ values produced unacceptable chattering in preliminary trials.
The cart-velocity damping term $-k_d\dot{x}$ is retained in the deployed
SMC+WO variant ($k_d{=}0.5$); without it, the post-balance cart-drift mode
described above remains open-loop and the cart eventually reaches a
rail. Earlier SMC formulations without the wall-override achieved only
about 12\% hardware SR.

\paragraph{Wall-proximity override.}
The same override as Eq.~\eqref{eq:override} is applied with
$d_w{=}0.12$~m, $u_\text{ov}{=}0.5$. Combined with the $k_d$ damping
term, this raises hardware SR from 12\% (the basic variant without
override or damping) to \textbf{92\%}.

\subsection{Linear MPC}

MPC formulates a finite-horizon QP using the linearized
first-order motor model with cart position as a hard constraint:
\begin{equation}
\begin{aligned}
  \min_{U}\ & \sum_{k=0}^{N-1} \left[ x_k^\top Q\, x_k + u_k^\top R\, u_k \right]
  + x_N^\top Q_f\, x_N \\
  \text{s.t. }\ & x_{k+1} = A_d x_k + B_d u_k,\quad x_0 = s_t,\\
  & |[x_k]_1| \leq x_\text{lim},\quad |u_k| \leq 0.9~\text{m/s}
\end{aligned}
\end{equation}
with $N{=}10$, solved via OSQP~\cite{osqp}.
The QP cart constraint is MPC's native boundary handling and yields
44\% hardware SR; the deployed wall-override variant
($d_w{=}0.12$, $u_\text{ov}{=}0.5$) drops the QP cart constraint and
relies on the override instead, raising SR to \textbf{80\%}.

\subsection{Nonlinear MPC (NMPC)}

NMPC solves a full nonlinear NLP using CasADi/IPOPT~\cite{casadi}
with RK4 discretization of the velocity-input dip dynamics:
\begin{equation}
\begin{aligned}
  \ddot\theta = \frac{1}{C_\mathrm{eff}(\theta)}\Big[
    &-R\cos\theta \cdot \ddot x - \tfrac{1}{2}C'(\theta)\dot\theta^2 \\
    &+ g(R\sin\theta - a(\theta)) - Q_\text{rr} - Q_\text{visc} \Big]
\end{aligned}
\end{equation}
with $a(\theta)$, $C_\mathrm{eff}(\theta)$, and the friction terms $Q_\text{rr}$,
$Q_\text{visc}$ as defined in Appendix~\ref{appendix:dynamics}, and cart
position enforced as hard NLP constraints over the $N{=}10$ horizon.

\paragraph{Beneficial model mismatch.}
NMPC's internal model uses FWHM$=$42~mm for the arc's dip width, wider
than the actual 20~mm. This wider-well approximation is beneficially
conservative and produces smoother recoveries, gentler commands that
avoid exciting the wall-bounce limit cycle
(Appendix~\ref{appendix:results}), yielding \textbf{90\% hardware SR}
(50 trials; first-order simulation predicts 100\%).

\subsection{Reinforcement Learning}

We train five algorithms in simulation using
StableBaselines3~\cite{stable_baselines3} and \texttt{sb3\_contrib}:
PPO, TRPO (on-policy, 16 parallel environments), and
SAC, TD3, TQC (off-policy, single environment).
PPO and TRPO use an MLP $[64,\,64]$ with \texttt{tanh} activations; the
off-policy algorithms SAC and TQC use the StableBaselines3 default
$[256,\,256]$, and TD3 its default $[400,\,300]$. Consistent with
the main paper, on-policy algorithms (PPO, TRPO) train for 3M
timesteps and off-policy algorithms (SAC, TD3, TQC) for 1M
timesteps, at 20~Hz.

\paragraph{Domain randomization (DR).}
At each episode reset, physical parameters are sampled uniformly
from ranges centered on measured values: gravity $\pm$0.5\%,
arc radius $\pm$5\%, cart friction $\pm$56\%, ball rolling
friction $\pm$60\%, ball viscous damping $\pm$67\%,
motor lag $\tau_v \pm$33\%, dip depth $\pm$50\%,
and dip width $\pm$40\%. Ranges are deliberately wide for
parameters that are difficult to measure precisely (friction,
damping) and narrow for well-characterized quantities (gravity,
arc radius). The nominal values and absolute ranges are gravity
$9.81\,[9.76, 9.86]$, arc radius $2.101\,[2.001, 2.201]$\,m, cart
friction $0.08\,[0.035, 0.125]$, ball rolling friction
$0.025\,[0.010, 0.040]$, ball viscous damping $0.009\,[0.003, 0.015]$,
motor lag $\tau_v\ 0.15\,[0.10, 0.20]$\,s, dip depth
$0.002\,[0.001, 0.003]$\,m, and dip width (FWHM)
$0.020\,[0.012, 0.028]$\,m, all specified in
\texttt{balancer/core/params.py}.

\paragraph{Blind-zone modeling (BZ).}
The simulation optionally clamps the $\theta$ observation to
$\pm\theta_\text{blind}$ and sets the velocity observation to $\approx 0$
when $|\theta| > \theta_\text{blind}$ (nominally 0.064\,rad from the sensor
geometry; the measured onset is ${\approx}0.071$\,rad) to reproduce the
VL53L0X minimum-range saturation
(Appendix~\ref{appendix:sensor_characterization}). The exact onset
$\theta_\text{blind}$ depends on the per-sensor calibration constants and
physical mounting, so under domain randomization the threshold is
sampled over $[0.060, 0.080]$\,rad rather than fixed. Combined with DR,
this yields the best RL performance:
PPO~(BZ+DR) achieves \textbf{96\% hardware SR}.

% -- Parameter Table ------------------------------------------
\begin{table}[htbp]
\centering
\renewcommand{\arraystretch}{1.15}
\caption{Deployed Controller Parameters}
\label{tab:controller_params}
\scriptsize
\setlength{\tabcolsep}{2pt}
\begin{tabular}{llc}
\toprule
\textbf{Controller} & \textbf{Parameter} & \textbf{Value} \\
\midrule
PD+WO   & $K_p,\; K_d$              & 19.78,\; 1.37   \\
        & $d_w,\; u_\text{ov}$      & 0.12~m,\; 0.5   \\
\midrule
LQR+WO  & $K$                       & $[0,\;0,\;{-19.78},\;{-1.37}]$ \\
        & $d_w,\; u_\text{ov}$      & 0.12~m,\; 0.5   \\
\midrule
SMC+WO  & $\lambda,\; k,\; \phi$    & 3.0,\; 6.0,\; 0.806 \\
        & $k_d$                     & 0.5              \\
        & $d_w,\; u_\text{ov}$      & 0.12~m,\; 0.5   \\
\midrule
MPC+WO  & $Q,\; R,\; N$             & diag(0.001,0,40,0),\; 0.017,\; 10 \\
        & $d_w,\; u_\text{ov}$      & 0.12~m,\; 0.5   \\
\midrule
NMPC    & $Q,\; R,\; N$             & diag(0,0,9,0),\; 0.006,\; 10 \\
        & FWHM (internal)           & 42~mm            \\
\midrule
RL      & Architecture              & MLP [64,64] on-pol.; [256,256] off-pol. \\
(online)& $\gamma,\; \lambda_\text{GAE}$ (on-pol.) & 0.99,\; 0.95 \\
        & Training steps            & on-policy 3M; off-policy 1M  \\
\midrule
IQL     & Architecture              & MLP [256, 256], ReLU \\
(offline)& Expectile, $\beta$        & 0.7,\; 3.0 \\
        & Training data             & $\approx$333k real transitions (post-recalib) \\
        & Gradient steps            & 280k (of 500k) \\
\midrule
MPPI    & Horizon $H$, samples $N$ & 30,\; 800 \\
        & $Q$ cost                  & diag($0,0.1,30,0.1$) \\
        & $R$, $\lambda$            & 0.001,\; 0.4 \\
        & World model               & 2-layer LSTM, see App.~\ref{appendix:world_model} \\
        & Runtime device            & CPU (faster than GPU for this small LSTM) \\
        & PyTorch threads           & 1 (conservative; 4 threads gives 20.07\,ms) \\
\bottomrule
\end{tabular}
\end{table}

\subsection{Offline Reinforcement Learning (IQL)}
\label{sec:iql_impl}

In addition to simulation-trained online RL, we evaluate
\emph{Implicit Q-Learning} (IQL)~\cite{kostrikov2022iql},
an offline RL algorithm trained entirely on real hardware
transitions from multiple controller families with no simulator access
during training.

\paragraph{Algorithm.}
IQL avoids querying out-of-distribution actions by learning a
state-value function $V(s)$ via expectile regression and
extracting a policy through advantage-weighted regression:
\begin{equation}
  \pi(a|s) \propto \exp\!\bigl(\beta\,(Q(s,a) - V(s))\bigr)
  \label{eq:iql_policy}
\end{equation}
where $\beta$ is the inverse temperature (\texttt{weight\_temp}).
This avoids the need for pessimistic Q-value penalties (as in CQL)
or a separate additive behavior-cloning penalty (as in TD3+BC), making IQL
well-suited to narrow mixed offline datasets.

\paragraph{Training data and hyperparameters.}
IQL is trained on $\approx$333k real post-recalibration transitions
(1{,}415 mixed-controller trials, no simulation data) with
d3rlpy~2.8.1~\cite{seno2022d3rlpy}. The dataset composition, network
sizes, and the full hyperparameter set (expectile 0.7,
$\beta = 3.0$, and the 280k-step checkpoint selection) are given in
Appendix~\ref{appendix:offline_rl}.

\paragraph{Deployment.}
The trained d3rlpy model is wrapped as an
\texttt{OfflineRLController} that accepts the standard
4-dimensional state and returns $u \in [-1,1]$ via deterministic
policy inference (no action noise at deployment).
